% Large Language Models (LLMs) have revolutionized reasoning and knowledge synthesis, yet remain static after training, unable to grow with experience as intelligent beings do. 
% This rigidity fundamentally constrains their adaptability in scientific domains that demand continual learning and evolving expertise.
% We propose \textit{\method{}}, a groundbreaking learning paradigm that enables LLM agents to continuously evolve without gradient updates.
% Instead of backpropagation, \method{} empowers agents to engage in expansive experiential exploration, reflecting on their reasoning trajectories to distill lessons of success and failure. 
% These distilled insights are structured into a self-evolving experience library, dynamically refined by semantic feedback signals that act as non-gradient supervisory forces.
% Instead of backpropagation, \method{} drives agents to conduct expansive experiential exploration, transforming their lessons of success and failure into a structured and continually evolving experience library.
% We propose \textit{\method{}}, a groundbreaking learning paradigm to drive agents to conduct expansive experiential exploration, transforming their lessons of success and failure into a structured and continually evolving experience library.
% During inference, agents retrieve and synthesize pertinent experiential knowledge, invoking learned cognitive patterns to tackle novel and complex tasks with greater reasoning adaptability and depth.
% The structured experience library serves as a transferable cognitive substrate, allowing other agents to inherit and apply accumulated intelligence in a plug-and-play way.
%Extensive experiments are conducted across diverse domains, including mathematical reasoning, chemical retro-synthesis, and protein fitness prediction, \method{} achieves up to 58\% improvement on AIME25 and 50\% on USPTO50k, outperforming strong LLM baselines.
% \method{} demonstrates substantial and consistent improvements on mathematical reasoning, chemical retro-synthesis, and protein fitness prediction (up to 58\% improvement on AIME25 and 50\% on USPTO50k).
% Intriguingly, we observe a distinct scaling law, where accumulated experiences lead to predictable and systematic performance gains as the experience library grows. 
% Beyond scaling, the experience library demonstrates intelligence inheritance, the ability to share experiential knowledge across different agents, resembling human-like cognitive transfer.
% Intriguingly, we observes a scaling law of experience growing and demonstrate knowledge transferability across different agents.
% Together, these results position \method{} as a human-analogous, gradient-free learning paradigm, illuminating a pathway toward self-evolving artificial intelligence.

Autonomous agents driven by Large Language Models (LLMs) have revolutionized reasoning and problem-solving but remain static after training, unable to grow with experience as intelligent beings do during deployment.
We introduce \textbf{F}orward \textbf{L}earning with \textbf{EX}perience~(\method{}), a gradient-free learning paradigm that enables LLM agents to continuously evolve through accumulated experience.
Specifically, \method{} cultivates scalable and inheritable evolution by constructing a structured experience library through continual reflection on successes and failures during interaction with the environment.
\method{} delivers substantial improvements on mathematical reasoning, chemical retrosynthesis, and protein fitness prediction (up to 23\% on AIME25, 10\% on USPTO50k, and 14\% on ProteinGym).
%(from 40\% to 63\% on AIME25, 20\% to 30\% on USPTO50k, and 46\% to 60\% on ProteinGym).
We further identify a clear scaling law of experiential growth and the phenomenon of experience inheritance across agents, marking a step toward scalable and inheritable continuous agent evolution.
Project Page: \url{https://flex-gensi-thuair.github.io}.